\documentclass[11pt,a4paper]{article}

\usepackage[utf8]{inputenc}
\usepackage[T1]{fontenc}
\usepackage[czech]{babel}
\usepackage[margin=2.7cm]{geometry}
\usepackage{amsmath,amssymb,amsthm,mathtools}
\usepackage{booktabs}
\usepackage{enumitem}
\usepackage{hyperref}
\usepackage{algorithm}
\usepackage{algpseudocode}
\usepackage{xcolor}
\usepackage{mathrsfs}

\newtheorem{lemma}{Lemma}
\newtheorem{theorem}{Věta}
\newtheorem{definition}{Definice}
\newtheorem{corollary}{Důsledek}
\newtheorem{remark}{Poznámka}

\newcommand{\R}{\mathbb{R}}
\newcommand{\code}[1]{\texttt{#1}}
\newcommand{\rank}{\operatorname{rank}}
\newcommand{\dimspan}{\operatorname{dim}}
\newcommand{\Var}{\operatorname{Var}}
\newcommand{\Cov}{\operatorname{Cov}}

\title{Systémová spolehlivost staticky neurčitých příhrad\\
\large Kompletní odvození: Fáze A -- C\\
\normalsize (null-space cut-sets podle Wei \& Deng 2022 + Cornell index / equivalent planes / PNET podle Rodrigues da Silva et al. 2024)}
\author{}
\date{}

\begin{document}
\maketitle
\tableofcontents

\bigskip
\noindent\textbf{Poznámka k tomuto dokumentu.} Cílem je odvodit \emph{proč} algoritmus implementovaný
v \code{system-reliability-truss-matlab/src/} funguje, ne jen zopakovat vzorce z pramenů. Kde je to
možné, je uveden i důkaz. Odkazy na konkrétní soubory/řádky (\code{soubor.m:řádek}) ukazují, kde přesně
je daný krok v kódu implementován, aby šlo teorii a implementaci propojit 1:1.

\section{Motivace a rozdíl oproti staticky určité příhradě}

U staticky určité příhrady (modul \code{reliability-truss-matlab}) platí: síla v každém prutu je
jednoznačně určena rovnováhou, takže selhání \emph{jediného} prutu = ztráta rovnováhy = kolaps celé
konstrukce. Systém je tedy \textbf{sériový}:
\[
  g_{\mathrm{sys}}(\mathbf X) = \min_{p=1,\dots,n_{\mathrm{el}}} g_p(\mathbf X).
\]

U staticky neurčité příhrady existuje $g_H>0$ nadbytečných prutů (stavů samo-napětí). Selhání jednoho
prutu vede k \textbf{redistribuci} vnitřních sil na zbylou konstrukci -- ta nemusí zkolabovat. Ke kolapsu
(mechanismu) dojde, až se odstraní dostatečně velká \emph{množina} prutů. Systémová spolehlivost proto
vyžaduje:

\begin{enumerate}[label=(\Alph*)]
  \item najít \textbf{všechny minimální množiny prutů} (tzv. \emph{cut-sets}), jejichž současné odstranění
        způsobí mechanismus (a žádná jejich vlastní podmnožina mechanismus nezpůsobí),
  \item pro každý cut-set spočítat, jak se síly \textbf{přerozdělují}, když pruty selhávají postupně
        v nějakém pořadí (\emph{failure sequence}),
  \item spočítat pravděpodobnost, že \emph{celý} cut-set selže (paralelní systém), tyto pravděpodobnosti
        přes všechny cut-sety zkombinovat do jednoho systémového čísla $\beta_{\mathrm{sys}}$.
\end{enumerate}

Tomu přesně odpovídají Fáze A, B, C níže.

\section{Fáze A -- Rovnovážná matice, null space a minimální cut-sets}
\label{sec:faze-a}

Zdrojové soubory: \code{equilibriumMatrixFn.m}, \code{nullSpaceCutSetsFn.m}.

\subsection{Odvození rovnovážné matice $A$}

Uvažujme 2D pin-jointed příhradu s $N$ volnými (neuvázanými) stupni volnosti uzlů a $n_{\mathrm{el}}$ pruty.
Nechť $\mathbf d\in\R^N$ je vektor uzlových posunů (volné DOFy), $\mathbf t\in\R^{n_{\mathrm{el}}}$ vektor
osových sil v prutech (kladně = tah).

\paragraph{Kompatibilita.} Prodloužení prutu $k$ mezi uzly $h$ (hlava) a $e$ (pata) je projekce relativního
posunu do osy prutu:
\[
  e_k = \mathbf u_k^\top(\mathbf d_e - \mathbf d_h), \qquad
  \mathbf u_k = \frac{1}{L_k}(\Delta x_k, \Delta z_k) \ \text{(jednotkový směrový vektor hlava→pata)}.
\]
Zapsáno přes globální DOF vektor $\mathbf d$: $e_k = \mathbf a_k^\top \mathbf d$, kde $\mathbf a_k\in\R^N$ má
nenulové složky jen ve 4 řádcích odpovídajících uzlům $h,e$ (v kódu: \code{codeNumbersFn}). Sestavením
$A := [\mathbf a_1,\dots,\mathbf a_{n_{\mathrm{el}}}]\in\R^{N\times n_{\mathrm{el}}}$ dostaneme kompatibilní
vztah pro všechny pruty najednou:
\[
  \mathbf e = A^\top \mathbf d.
\]

\paragraph{Rovnováha (kontragredience).} Princip virtuálních prací: vnější virtuální práce = vnitřní
virtuální práce pro libovolné $\delta\mathbf d$:
\[
  \mathbf f^\top \delta\mathbf d \;=\; \sum_k t_k\,\delta e_k \;=\; \sum_k t_k\,(\mathbf a_k^\top\delta\mathbf d)
  \;=\; \Big(\sum_k t_k \mathbf a_k\Big)^{\!\top}\delta\mathbf d \;=\; (A\mathbf t)^\top \delta\mathbf d.
\]
Protože to platí pro \emph{libovolné} $\delta \mathbf d$, plyne $\mathbf f = A\mathbf t$ -- tedy \textbf{stejná}
matice $A$, která popisuje kompatibilitu ($\mathbf e=A^\top\mathbf d$), popisuje i rovnováhu
($\mathbf f = A\mathbf t$). Toto je klasický výsledek teorie tuhosti/mechanismů (Pellegrino \& Calladine
1986) a je důvodem, proč se $A$ říká \emph{rovnovážná matice} (equilibrium matrix), přestože vznikla
z kompatibility. Znaménková konvence jednotlivých sloupců $\mathbf a_k$ je věc volby (v kódu
\code{equilibriumMatrixFn.m:42-54} je hlava $+[c,s]$, pata $-[c,s]$) -- simultánní obrácení znaménka
celého sloupce nemění $\rank(A)$ ani jeho nulový prostor, což je jediné, na čem dále záleží.

\subsection{Dvě fundamentální podmnožiny -- stavy samo-napětí a mechanismy}

Z rozkladu $\R^{n_{\mathrm{el}}} = \operatorname{null}(A)\oplus(\text{ostatní})$ a věty o dimenzi
jádra a obrazu (rank--nullity):

\begin{itemize}
  \item $\operatorname{null}(A) = \{\mathbf s : A\mathbf s = 0\}$ -- vektory vnitřních sil, které jsou
        v rovnováze \textbf{bez jakéhokoli vnějšího zatížení} -- tzv. \textbf{stavy samo-napětí}
        (self-stress states). Jejich počet je
        \[
          g_H := \dim\operatorname{null}(A) = n_{\mathrm{el}} - \rank(A)
        \]
        -- to je přesně \textbf{stupeň statické neurčitosti}.
  \item $\operatorname{null}(A^\top) = \{\mathbf d: A^\top\mathbf d=0\}$ -- posuny, které nezpůsobí
        prodloužení \emph{žádného} prutu -- tj. \textbf{mechanismy}. Protože jsme už odečetli podepřené
        DOFy (žádné volné tuhé pohyby), platí: konstrukce je \textbf{kinematicky stabilní} (bez
        mechanismu) $\iff \operatorname{null}(A^\top)=\{0\} \iff \rank(A) = N$.
\end{itemize}

V dalším textu předpokládáme (Fáze A pracuje s neporušenou konstrukcí), že $\rank(A)=N$, tj. konstrukce
je stabilní a $g_H = n_{\mathrm{el}}-N \geq 0$ je počet nadbytečných prutů. Necht' $V\in\R^{n_{\mathrm{el}}
\times g_H}$ je ortonormální báze $\operatorname{null}(A)$ (v kódu získaná ze SVD:
\code{nullSpaceCutSetsFn.m:97-104}, $AV=0$).

\subsection{Lemma 1 -- nutný prut (\emph{necessary bar}, cut-set velikosti 1)}

\begin{lemma}[Necessary bar]
Prut $i$ je \emph{nutný} (jeho samotné odstranění způsobí mechanismus, tj. $\rank(A_{-i}) = N-1$)
právě tehdy, když $i$-tý řádek $V_i$ matice $V$ je nulový, tj. $s_i = 0$ pro \textbf{každý} stav
samo-napětí $\mathbf s\in\operatorname{null}(A)$.
\end{lemma}

\begin{proof}
Odstranění prutu $i$ z matice $A$ = odstranění sloupce $\mathbf a_i$, čímž vznikne
$A_{-i}\in\R^{N\times(n_{\mathrm{el}}-1)}$. Protože se odstraňuje jeden sloupec, $\rank(A_{-i})\in\{N-1,N\}$.

\smallskip
\noindent($\Leftarrow$, existuje $\mathbf s$ se $s_i\neq0$ $\Rightarrow$ prut NENÍ nutný): Z
$A\mathbf s=0$ plyne $\sum_k s_k\mathbf a_k = 0$, tedy
\[
  \mathbf a_i = -\frac1{s_i}\sum_{k\neq i} s_k\, \mathbf a_k \ \in\ \operatorname{span}\{\mathbf a_k:k\neq i\}.
\]
Přidání $\mathbf a_i$ tedy nerozšíří rozpětí sloupců, takže $\operatorname{span}(A)=\operatorname{span}
(A_{-i})$ a $\rank(A_{-i})=\rank(A)=N$ -- žádný mechanismus.

\smallskip
\noindent($\Rightarrow$, kontrapozitivně: pokud $\rank(A_{-i})=N$, pak existuje $\mathbf s$ se
$s_i\neq0$): Je-li $\rank(A_{-i})=N$, pak $\mathbf a_i\in\operatorname{span}\{\mathbf a_k:k\neq i\}$,
tj. existují $c_k$ takové, že $\mathbf a_i = \sum_{k\neq i} c_k \mathbf a_k$. Položme
$s_i:=1,\ s_k:=-c_k\ (k\neq i)$. Pak
\[
  \sum_k s_k \mathbf a_k = \mathbf a_i - \sum_{k\neq i} c_k\mathbf a_k = 0,
\]
takže $\mathbf s\in\operatorname{null}(A)$ a $s_i=1\neq0$ -- QED (kontrapozice dává původní tvrzení).
\end{proof}

Toto přesně odpovídá kódu \code{nullSpaceCutSetsFn.m:113-125}: řádky $V$ s normou $\approx0$
$\Rightarrow$ cut-set velikosti 1.

\subsection{Lemma 2 -- obecný cut-set (zobecnění Lemma 1)}

\begin{lemma}[Cut-set kritérium]
Necht' $Q\subseteq\{1,\dots,n_{\mathrm{el}}\}$, $|Q|=k$, a $V_Q\in\R^{k\times g_H}$ je submatice $V$
tvořená řádky $Q$. Pak
\[
  \rank(A_{-Q}) \;=\; N - \big(k - \rank(V_Q)\big).
\]
Speciálně, odstranění $Q$ způsobí mechanismus $\iff \rank(V_Q) < k$.
\end{lemma}

\begin{proof}
Označme $e_{Q}\subset\R^{n_{\mathrm{el}}}$ souřadnicový podprostor odpovídající indexům $Q$ (dimenze
$k$) a $e_{Q^c}$ jeho komplement (dimenze $n_{\mathrm{el}}-k$). Sloupcový prostor $A_{-Q}$ je právě
obraz $A(e_{Q^c})$ (obraz vektorů, které mají nulové složky v $Q$). Podle věty o dimenzi jádra a obrazu
aplikované na zúžené zobrazení $A|_{e_{Q^c}}$:
\[
  \dim A(e_{Q^c}) \;=\; \dim e_{Q^c} \;-\; \dim\big(e_{Q^c}\cap \operatorname{null}(A)\big)
  \;=\; (n_{\mathrm{el}}-k) - \dim\big(\operatorname{null}(A)\cap e_{Q^c}\big).
\]
Nyní $\operatorname{null}(A)\cap e_{Q^c} = \{\mathbf s = V\mathbf c : V_Q\mathbf c = 0\}$ (stavy
samo-napětí, které mají nulovou sílu ve VŠECH prutech $Q$) $=V\cdot\operatorname{null}(V_Q)$, což má
dimenzi $\dim\operatorname{null}(V_Q) = g_H - \rank(V_Q)$ (rank--nullity pro $V_Q:\R^{g_H}\to\R^k$;
$V$ má plnou sloupcovou hodnost $g_H$, takže $V\cdot(\cdot)$ zachovává dimenzi).

Dosazením a s využitím $n_{\mathrm{el}} - g_H = N$:
\[
  \rank(A_{-Q}) = (n_{\mathrm{el}}-k) - (g_H-\rank(V_Q)) = (n_{\mathrm{el}}-g_H) - k + \rank(V_Q)
  = N - k + \rank(V_Q).
\]
Mechanismus nastane $\iff \rank(A_{-Q}) < N \iff \rank(V_Q) < k$. \qed
\end{proof}

\begin{remark}
Pro $k=1$ dostáváme přesně Lemma 1 zpět ($\rank(V_{\{i\}})\in\{0,1\}$, $<1 \iff$ nulový řádek).
\end{remark}

\begin{corollary}[Maximální velikost cut-setu]
Protože $V_Q$ má nejvýše $g_H$ sloupců, platí vždy $\rank(V_Q)\le g_H$. Pro $k=g_H+1$ je tedy
$\rank(V_Q)\le g_H < k$ \textbf{triviálně}, nezávisle na geometrii -- libovolná $(g_H{+}1)$-tice prutů
je (dependentní) kandidát na cut-set. Žádný cut-set nemůže mít víc než $g_H+1$ prvků.
\end{corollary}

\subsection{Minimalita}

$Q$ (s $\rank(V_Q)<k$) je \textbf{minimální} cut-set, právě když žádná jeho vlastní podmnožina už
mechanismus nezpůsobuje, tj. \textbf{každá} $(k{-}1)$-podmnožina $Q'\subset Q$ splňuje
$\rank(V_{Q'}) = k-1$ (plná hodnost). Je-li totiž nějaká $(k{-}1)$-podmnožina sama dependentní, byla by
už (menším) cut-setem, a $Q$ by ho jen zbytečně obsahoval navíc -- není tedy minimální.

\subsection{Algoritmus (\code{nullSpaceCutSetsFn.m})}

\begin{algorithm}[h]
\caption{Nalezení všech minimálních cut-sets}
\begin{algorithmic}[1]
\State Sestav $A$ (\code{equilibriumMatrixFn}); SVD $\Rightarrow$ $\rank(A)$, $g_H = n_{\mathrm{el}}-\rank(A)$, báze $V$ jádra
\If{$g_H = 0$} \Return $\emptyset$ (staticky určitá, žádná redundance) \EndIf
\State Řádky $V$ s normou $\approx 0$ $\Rightarrow$ cut-sets $\{i\}$ (Lemma 1); odstraň tyto řádky $\to V'$ ($b'\times g_H$)
\For{$k = 2,\dots,g_H+1$}
    \For{každou $k$-podmnožinu $Q$ z $b'$ zbylých řádků}
        \If{$\rank(V'_Q) < k$} \Comment{Lemma 2 -- dependentní}
            \If{každá $(k{-}1)$-podmnožina $Q$ má plnou hodnost $k{-}1$} \Comment{minimalita}
                \State přidej $Q$ (přepočtenou zpět na původní čísla prutů) do seznamu cut-setů
            \EndIf
        \EndIf
    \EndFor
\EndFor
\State \Return seznam cut-setů (deduplikován jako množiny)
\end{algorithmic}
\end{algorithm}

Cena: pro $k=2,\dots,g_H{+}1$ testujeme $\binom{b'}{k}$ podmnožin -- proveditelné, protože $g_H$ bývá
malé ($\le\!\sim\!7$ i pro reálné konstrukce, viz \code{metodika.md}).

\subsection{Zjištěná mezera v algoritmu Wei \& Deng (2022) pro $k=g_H+1$}

Wei \& Deng (2022) ve své vlastní kapitole 4 zaznamenávají cut-set jen tehdy, pokud koeficienty
$\gamma_p = (V_b^\top)^{-1}\mathbf v_p$ (rozklad sloupce $p$ do zvolené báze $g_H$ sloupců) obsahují
\textbf{alespoň jednu nulu}. To ale strukturálně nemůže zachytit cut-sety velikosti přesně $g_H+1$,
u kterých je -- podle Corollary výše -- dependence \textbf{triviální} (plyne z počtu stupňů volnosti, ne
z konkrétních hodnot $\gamma_p$), takže obecně \textbf{žádný} $\gamma_p$ v žádné použitelné bázi nemusí
mít nulovou složku, i když jde o platný minimální cut-set.

\textbf{Konkrétní ověřený příklad} (15-prutová příhrada z Wei \& Deng, kap. 6.1, geometrie ověřena
pixel-přesně z jejich Fig. 3): trojice $\{1,3,8\}$ je stabilní, staticky určitá podmnožina
($\rank(V_{\{1,3,8\}})=3=g_H$, tedy platná báze $V_b$). Přidáním libovolného dalšího prutu z báze
disjunktní čtveřice, např. prutu 10, vznikne $\{1,3,8,10\}$ -- matematicky nutně mechanismus
($\rank(V_Q)\le g_H=3 < 4=k$), ověřeno i přímým výpočtem $\rank(A_{-Q})$ na plné matici a robustností
vůči perturbaci geometrie o $0{,}01\,\%$. Wei \& Deng ve své Tabulce 1 pro $p=10$ s bází $\{1,3,8\}$
uvádí $\gamma=[-1,1,1]$ (žádná nula) $\Rightarrow$ nezaznamenáno -- ačkoli jde o platný minimální
cut-set. Náš kompletní (Lemma 1 + Lemma 2, bez zkratek algoritmu) postup najde touto cestou navíc
\textbf{25} cut-setů velikosti 4 (tvaru $\{$jeden z $[1,2,5,11,12]\}\times\{$jeden z
$[3,4,7,14,15]\}\times\{8,10\}$), celkem $32+25=57$ oproti publikovaným 32. Numericky na výsledné
$\beta_{\mathrm{sys}}$ to nemá prakticky vliv (cut-sety velikosti $g_H{+}1$ mají vyšší $\beta_p$, tedy
zanedbatelný příspěvek a navíc vysokou korelaci s menšími podmnožinami, které PNET stejně pohltí) --
implementace v \code{nullSpaceCutSetsFn.m} proto úmyslně \emph{nereplikuje} tuto mezeru a hledá úplně.

\section{Fáze B -- Vlivové koeficienty na redukované konstrukci}
\label{sec:faze-b}

Zdrojový soubor: \code{reducedStructureInfluenceFn.m}.

Fáze B je čistě FEM (žádné náhodné veličiny) a poskytuje "stavební kameny" pro linearizovanou mezní
funkci ve Fázi C. Pro danou (postupně redukovanou -- tj. s odstraněnými pruty) konstrukci potřebujeme
dvě věci:

\subsection{Koeficient $b_{ij}$ -- vliv vnějšího zatížení}

$b_{ij}$ := osová síla v prutu $i$ vyvolaná \textbf{jednotkovým} zatížením typu $j$ (jednotková síla /
kombinace) aplikovaným na konstrukci, ze které jsou už odstraněny pruty v $Q\ (=\code{removedMembers})$.
Získá se přímo řešením redukované soustavy tuhosti:
\[
  K_{Q}\,\mathbf d = \mathbf f_j \quad(\text{jednotkové zatížení } j) \;\Longrightarrow\;
  b_{ij} = N_i(\mathbf d) \quad \text{(zpětný výpočet osové síly v prutu $i$)}.
\]
Protože FEM je lineární, skutečná síla v prutu $i$ od zatížení $\mathbf P=(P_1,\dots,P_{n_L})$ na TÉŽE
redukované konstrukci je pak přímo superpozicí:
\[
  N_i \;=\; \sum_j b_{ij}\,P_j.
\]
V kódu: \code{reducedStructureInfluenceFn.m:93-100} (\code{linearSolverFn} na \code{membersRed}, sloupec
1 z \code{endForces.local}).

\subsection{Koeficient $a_{il}$ -- vliv zbytkové síly z redistribuce}

Když prut $l$ selže (křehce nebo tažně), do okamžiku selhání nesl sílu $\eta_l R_l$ ($\eta_l\in[0,1]$:
$0$=křehké/žádná zbytková síla, $1$=plně tažné/celá síla se přenese). Tato síla se v okamžiku selhání
"uvolní" jako pár samo-vyrovnaných sil na koncových uzlech prutu $l$ (ve směru jeho původní osy) a
zbytek konstrukce ji musí převzít. $a_{il}$ := osová síla v prutu $i$ od \textbf{jednotkové} takové
samo-vyrovnané dvojice sil, na konstrukci, ze které je prut $l$ (a případně další už selhané pruty)
odstraněn:
\[
  \mathbf f_l = (+u_x,+u_z)\big|_{\text{uzel hlava}(l)} \ \oplus\ (-u_x,-u_z)\big|_{\text{uzel pata}(l)}
  \qquad (u_x,u_z)=\text{směrové kosiny prutu } l.
\]
V kódu: \code{reducedStructureInfluenceFn.m:110-128}. Skutečná zbytková síla je $\eta_l R_l$, takže
příspěvek do celkové síly v prutu $i$ je opět superpozicí: $a_{il}\cdot(\eta_l R_l)$.

\subsection{Poznámka k singularitě}

Jakmile odstraněná množina prutů obsahuje celý cut-set, $K_Q$ je singulární (mechanismus) --
\code{linearSolverFn} vydá varování o RCOND. Toto varování je \emph{spolehlivý} signál mechanismu
(viz \code{metodika.md} kap. 6); samotná konečnost výsledku \code{b\_ij}/\code{a\_il} se na to
nedá spolehnout (zpětná substituce na singulární matici stále vrátí nějaké konečné, ale fyzikálně
nesmyslné číslo).

\section{Fáze C -- Limitní funkce, Cornellův index, ekvivalentní roviny, PNET}
\label{sec:faze-c}

Zdrojové soubory: \code{sequenceLimitStateFn.m}, \code{cornellIndexFn.m},
\code{mostProbableSequenceFn.m}, \code{equivalentPlaneFn.m}, \code{pnetSystemReliabilityFn.m},
\code{systemReliabilityFn.m}.

\subsection{Linearizovaná mezní funkce prutu $i$ na pozici sekvence}
\label{sec:c1}

Uvažujme konkrétní \textbf{uspořádanou} posloupnost selhávání (permutaci) cut-setu,
$(1),(2),\dots,(m)$ (v kódu \code{sequence}), a soustřeďme se na prvek na pozici $\mathrm{pos}$,
$i=\mathrm{sequence}(\mathrm{pos})$, poté co už selhaly pruty na pozicích $1,\dots,\mathrm{pos}{-}1$.

\paragraph{Odvození $g_i$ ze superpozice.} Na konstrukci redukované o pruty $1,\dots,\mathrm{pos}{-}1$
(Fáze B) je celková síla v ještě stojícím prutu $i$ (lineární superpozice vnějšího zatížení a
zbytkových sil z dřívějších selhání):
\[
  N_i \;=\; \sum_j b_{ij} P_j \;+\; \sum_{l<i} a_{il}\,\eta_l R_l.
\]
Prut $i$ přežije, dokud $|N_i| \le R_i$ (kapacita $R_i$); porušen je ve směru daném znaménkem $N_i$, tj.
podmínkou přežití je $\operatorname{sign}(N_i)\cdot N_i \le R_i$. Definujeme bezpečnostní rezervu (mezní
funkci) tak, aby $g_i>0$ znamenalo "přežil", $g_i\le0$ "selhal":
\[
  \boxed{\;g_i \;=\; R_i \;-\; \operatorname{sign}(N_i)\Big(\sum_j b_{ij}P_j + \sum_{l<i} a_{il}\,\eta_l R_l\Big)\;}
  \tag{C.1}
\]
což je přesně rov. 14 v metodika.md / rov. 36 u Wei \& Deng. Znaménko $\operatorname{sign}(N_i)$ se v
implementaci bere ze \textbf{středního} (očekávaného) kombinovaného zatížení na téže redukované
konstrukci (\code{sequenceLimitStateFn.m:112-116}) -- jednoznačná volba, když působí víc nezávislých
zatížení současně.

\paragraph{Standardizace na sdílený gaussovský vektor $U$.} Aby $g_i$ vyšlo \textbf{přesně lineární}
v nezávislých standardizovaných normálních veličinách (nutná podmínka pro Cornellův index níže),
modelujeme kapacitu i zatížení přímo jako gaussovské náhodné veličiny:
\[
  R_i = \mu_{R,g(i)} + \sigma_{R,g(i)}\,U_{g(i)}, \qquad P_j = \mu_{P,j} + \sigma_{P,j}\,U_{n_G+j},
\]
kde $g(i)\in\{1,\dots,n_G\}$ je index skupiny průřezu prutu $i$ (pruty stejného typu sdílí TUTÉŽ
náhodnou veličinu $R$, tedy $\rho=1$ uvnitř skupiny) a
$U=(U_1,\dots,U_{n_G},U_{n_G+1},\dots,U_{n_G+n_L})\sim\mathcal N(\mathbf 0,I)$ je jeden sdílený vektor
přes \textbf{celou} analýzu (\code{sequenceLimitStateFn.m:70-74}). Dosazením do (C.1) a sečtením podle
složek $U$ dostaneme $g_i$ jako lineární formu
\[
  g_i \;=\; \overline g_i \;+\; \mathbf c_i^\top U, \qquad \overline g_i:=\mathbb E[g_i],
\]
kde (srovnej \code{sequenceLimitStateFn.m:118-140}):
\[
\begin{aligned}
  \overline g_i &= \mu_{R,g(i)} - \operatorname{sign}(N_i)\Big[\sum_j b_{ij}\mu_{P,j}
        + \sum_{l<i} a_{il}\,\eta_l\,\mu_{R,g(l)}\Big], \\[4pt]
  c_i[\,g(i)\,] &\mathrel{+}= \sigma_{R,g(i)}, \\
  c_i[\,n_G+j\,] &\mathrel{-}= \operatorname{sign}(N_i)\cdot b_{ij}\cdot\sigma_{P,j} \qquad(\forall j),\\
  c_i[\,g(l)\,] &\mathrel{-}= \operatorname{sign}(N_i)\cdot a_{il}\,\eta_l\cdot\sigma_{R,g(l)}
        \qquad(\forall l<i,\ \eta_l\neq0).
\end{aligned}
\tag{C.2}
\]
Protože složky $U$ jsou nezávislé a jednotkového rozptylu,
\[
  \Var[g_i] \;=\; \|\mathbf c_i\|_2^2 \;=\; \sum_m c_i[m]^2. \tag{C.3}
\]
(\code{seq(pos).coeffU}, \code{.meanG}, \code{.varG} v kódu -- (C.2)--(C.3) jsou \code{sequenceLimitStateFn.m:120-140}
doslovně.)

\subsection{Cornellův index sekvence -- paralelní systém}
\label{sec:c2}

Sekvence (jako celek) selže, jen když selžou \textbf{všechny} její prvky současně (paralelní podsystém):
$\{g_1<0\}\cap\dots\cap\{g_m<0\}$. Standardizujme:
\[
  Z_i := \frac{g_i}{\mathrm{sd}_i},\qquad \mathrm{sd}_i := \sqrt{\Var[g_i]} = \|\mathbf c_i\|,
  \qquad \beta_i := \frac{\overline g_i}{\mathrm{sd}_i}.
\]
Protože $g_i$ je gaussovská (lineární kombinace gaussovských $U$), $Z_i\sim\mathcal N(\beta_i,1)$ a
\[
  P(g_i<0) = P(Z_i<0) = \Phi(-\beta_i).
\]

\paragraph{Korelace mezi prvky.} Definujme jednotkový vektor citlivosti $\alpha_i := \mathbf c_i/\mathrm{sd}_i$
(jednotkový, protože $\mathrm{sd}_i=\|\mathbf c_i\|$). Pak, protože $U$ má nezávislé jednotkové složky:
\[
  \Cov(Z_i,Z_j) = \frac{\Cov(g_i,g_j)}{\mathrm{sd}_i\,\mathrm{sd}_j}
  = \frac{\mathbf c_i^\top \mathbf c_j}{\mathrm{sd}_i\,\mathrm{sd}_j} = \alpha_i^\top \alpha_j =: \rho_{ij}.
\tag{C.4}
\]
(V implementaci \code{cornellIndexFn.m:47} se $\rho=\alpha\alpha^\top$ počítá přesně takto; kód navíc
bere $\alpha_i^{\mathrm{code}} = -\mathbf c_i/\mathrm{sd}_i$ -- opačné znaménko oproti mému $\alpha_i$ výše
-- což je jen globální konvence: protože $\rho_{ij}=\alpha_i^\top\alpha_j$ je bilineární, simultánní
obrácení znaménka VŠECH $\alpha_i$ dá stejné $\rho$, takže výsledek je totožný.)

\paragraph{Sdružená pravděpodobnost.} $Z=(Z_1,\dots,Z_m)^\top \sim \mathcal N(\beta,\rho)$
(vícerozměrně normální s vektorem středních hodnot $\beta$ a korelační maticí $\rho$ z (C.4)). Chceme
\[
  P_{f,p} := P\Big(\bigcap_i \{Z_i<0\}\Big) = P(Z\le \mathbf 0).
\]
Standardní vícerozměrná normální CDF je definována pro $Y\sim\mathcal N(\mathbf 0,\Sigma)$ jako
$\Phi_m(\mathbf x;\Sigma)=P(Y\le\mathbf x)$. Substitucí $Y:=Z-\beta\sim\mathcal N(\mathbf 0,\rho)$:
\[
  P_{f,p} = P(Z\le\mathbf 0) = P(Y\le -\beta) = \Phi_m(-\beta;\rho).
\tag{C.5}
\]
a systémový (sekvenční) index spolehlivosti sekvence
\[
  \boxed{\;\beta_p := -\Phi^{-1}\big(\Phi_m(-\beta;\rho)\big)\;}
\tag{C.6}
\]
-- (C.5)--(C.6) jsou přesně Wei \& Deng rov. 38--39 a odpovídají \code{cornellIndexFn.m:47-78}
(\code{mvncdf((-betaVec)', zeros(1,m), rho)}).

\begin{remark}[Kontrola pro $m=1$]
Pro jediný prvek $\Phi_m$ degeneruje na $\Phi$, takže $\beta_p = -\Phi^{-1}(\Phi(-\beta_1)) = \beta_1$
-- konzistentní s jednorozměrným případem (\code{cornellIndexFn.m:26-31}).
\end{remark}

\subsection{Nejpravděpodobnější sekvence -- reprezentant cut-setu}
\label{sec:c3}

Cut-set z Fáze A je \textbf{neuspořádaná} množina; $\beta_p$ z (C.6) ale závisí na pořadí (protože
$b_{ij},a_{il}$ z Fáze B jsou počítány na postupně redukované konstrukci, a $a_{il}$ existuje jen
pro $l$ selhavší dříve). \code{mostProbableSequenceFn.m} tedy vyzkouší \textbf{všechny} permutace
$\pi$ prutů cut-setu, pro každou spočte $\beta_p(\pi)$ podle (C.1)--(C.6), a vybere
\[
  \beta_p^{\ast} := \min_{\pi\in\mathrm{Perm}(\text{cut-set})} \beta_p(\pi)
\]
(nejnižší $\beta_p$ = nejpravděpodobnější, tedy "reprezentativní", pořadí selhávání). Proveditelné jen
díky tomu, že $|{\rm cut\text{-}set}|\le g_H+1$ je malé (max. $5!=120$ permutací pro $g_H\le4$ typické
konstrukce).

\subsection{Ekvivalentní lineární rovina (linearizace reprezentanta)}
\label{sec:c4}

Problém: $\beta_p$ z (C.6) NENÍ lineární funkce $U$ (vzniklo z $\Phi_m$), takže různé cut-sety
nelze mezi sebou přímo korelačně srovnávat (PNET níže potřebuje $\rho_{kl}=\tilde\alpha_k^\top
\tilde\alpha_l$ mezi cut-sety). Wei \& Deng navrhují nahradit reprezentanta cut-setu $k$ jedním
\textbf{ekvivalentním lineárním} výkonovým funkcionálem
\[
  \tilde Z_k = \tilde\alpha_k^\top U + \tilde\beta_k,
\]
určeným dvěma požadavky:
\begin{enumerate}
  \item stejná pravděpodobnost: $\tilde\beta_k := \beta_p^\ast$ (rov. 45),
  \item stejná \textbf{prvořádová citlivost} $\beta_p$ na malé posunutí středních hodnot $U$
        (rov. 46--47).
\end{enumerate}

\paragraph{Odvození citlivosti.} Uvažujme myšlený posun středních hodnot $U\to U+\varepsilon$
($\varepsilon\in\R^n$ malé). Protože $Z_i=\alpha_i^\top U+\beta_i$, posun o $\varepsilon$ posune
efektivní $\beta_i$ o $\alpha_i^\top\varepsilon$:
\[
  \beta_i(\varepsilon) = \beta_i + \alpha_i^\top \varepsilon.
\]
$\beta_p$ je funkcí $\beta$ (přes (C.6), s $\rho$ fixní -- $\rho$ na $\varepsilon$ nezávisí, protože
posun středních hodnot nemění korelaci). Definujme $\beta_p(\varepsilon) := \beta_p(\beta(\varepsilon))$
přes (C.6) a hledejme
\[
  \operatorname{grad}_j := \frac{\partial \beta_p(\varepsilon)}{\partial \varepsilon_j}\bigg|_{\varepsilon=0}
  = \sum_i \frac{\partial\beta_p}{\partial\beta_i}\cdot\frac{\partial\beta_i(\varepsilon)}{\partial\varepsilon_j}
  = \sum_i \frac{\partial\beta_p}{\partial\beta_i}\,\alpha_i(j).
\]
Analytická derivace $\partial\beta_p/\partial\beta_i$ přes $\Phi_m^{-1}$/mnohorozměrnou CDF nemá v MATLABu
uzavřený tvar (Genz algoritmus pro \code{mvncdf} nedává gradient), takže se počítá
\textbf{centrální diferencí} přesně podle parametrizace výše:
\[
  \operatorname{grad}_j \approx \frac{\beta_p(\beta + h\,\alpha_{:,j}) - \beta_p(\beta - h\,\alpha_{:,j})}{2h},
  \qquad h=10^{-3},
\tag{C.7}
\]
kde $\beta_p(\cdot)$ značí (C.6) vyhodnocené s posunutým vektorem $\beta$ ($\alpha_{:,j}$ = $j$-tý
sloupec matice $\alpha$, tj. citlivost všech $m$ prvků sekvence na $U_j$). Přesně toto počítá
\code{equivalentPlaneFn.m:53-63}. Nakonec normalizace na jednotkovou délku:
\[
  \tilde\alpha_k := \frac{\operatorname{grad}}{\|\operatorname{grad}\|}. \tag{C.8}
\]

\subsection{PNET -- agregace přes všechny cut-sety}
\label{sec:c5}

Máme $q$ cut-setů, každý reprezentovaný lineárním $(\tilde\beta_k,\tilde\alpha_k)_{k=1}^q$ z
(C.6)--(C.8). Chceme \emph{systémovou} pravděpodobnost poruchy
\[
  P_{f,\mathrm{sys}} = P\Big(\bigcup_{k=1}^q \{\tilde Z_k < 0\}\Big).
\]
Přímý výpočet $q$-rozměrného sdruženého integrálu (typicky $q\sim10^1$--$10^2$) je numericky
neúnosný/nestabilní. PNET (Probabilistic NETwork) tento problém obchází pomocí korelace mezi cut-sety,
$\rho_{kl} = \tilde\alpha_k^\top\tilde\alpha_l$ (stejná bilineární forma jako (C.4), teď mezi
reprezentanty cut-setů):

\begin{algorithm}[h]
\caption{PNET seskupení a kombinace}
\begin{algorithmic}[1]
\State Seřaď cut-sety vzestupně podle $\tilde\beta_k$ (nejnižší = nejnebezpečnější první)
\While{existuje neseskupený cut-set}
  \State $r \gets$ aktuálně nejhorší (nejnižší $\tilde\beta$) neseskupený cut-set (nový reprezentant)
  \State spočti $\rho_{r,k} = \tilde\alpha_r^\top\tilde\alpha_k$ pro všechny neseskupené $k$
  \State všechny $k$ s $\rho_{r,k}\ge \rho_0$ (typicky $0{,}7$--$0{,}85$) přiřaď do skupiny reprezentanta $r$
\EndWhile
\State $w$ skupin, každá s reprezentativním $\beta_i := \tilde\beta_{r_i}$
\State $P_{f,\mathrm{sys}} = 1-\prod_{i=1}^w(1-\Phi(-\beta_i))$; \quad $\beta_{\mathrm{sys}}=-\Phi^{-1}(P_{f,\mathrm{sys}})$
\end{algorithmic}
\end{algorithm}

\paragraph{Zdůvodnění kroku 8 (nezávislý sériový systém).} Cut-sety uvnitř jedné skupiny jsou (skoro)
\textbf{plně korelované} ($\rho\ge\rho_0$) s reprezentantem -- prakticky "stejná" informace o poruše,
tedy se počítají jen jednou (přes reprezentanta). Mezi RŮZNÝMI skupinami je korelace $<\rho_0$ (podle
konstrukce), tedy se aproximují jako \textbf{nezávislé}. Pro nezávislé jevy $E_1,\dots,E_w$ platí
klasicky
\[
  P(\text{přežití}) = P\Big(\bigcap_i \overline{E_i}\Big) = \prod_i P(\overline{E_i}) = \prod_i(1-P(E_i))
  \;\Longrightarrow\; P_{f,\mathrm{sys}} = 1-\prod_i(1-P_{f,i}),
\tag{C.9}
\]
což je přesně krok 8 (Wei \& Deng rov. 48) a \code{pnetSystemReliabilityFn.m:65-67}. Toto je
\textbf{přesné} pro skutečně nezávislé jevy a \textbf{aproximace} v tom rozsahu, v jakém
skupiny nejsou zcela nezávislé (řízeno prahem $\rho_0$ -- kompromis mezi přesností a výpočetní
náročností plného $q$-rozměrného integrálu).

\subsection{Orchestrace celého pipeline (\code{systemReliabilityFn.m})}

\begin{algorithm}[h]
\caption{Systémová spolehlivost staticky neurčité příhrady -- celý postup A--C}
\begin{algorithmic}[1]
\State \textbf{Fáze A}: $A\gets$\code{equilibriumMatrixFn}; $(\text{cutSets}, g_H)\gets$\code{nullSpaceCutSetsFn}$(A)$
\For{každý cut-set $c=1,\dots,q$}
    \State \textbf{Fáze C.3}: $(\pi^\ast,\beta_p^\ast) \gets$ \code{mostProbableSequenceFn} (brute-force přes permutace, uvnitř volá \textbf{Fáze B} + \textbf{C.1--C.2})
    \State \textbf{Fáze C.4}: $(\tilde\beta_c,\tilde\alpha_c)\gets$\code{equivalentPlaneFn}$(\alpha,\beta\,|\,\pi^\ast)$
\EndFor
\State \textbf{Fáze C.5}: $(\beta_{\mathrm{sys}}, P_{f,\mathrm{sys}}, \text{skupiny})\gets$\code{pnetSystemReliabilityFn}$(\tilde\beta,\tilde\alpha,\rho_0)$
\State \Return $\beta_{\mathrm{sys}}$ + diagnostika (\code{betaPerCutSet}, \code{pnetGroups}, ...)
\end{algorithmic}
\end{algorithm}

\section{Shrnutí a ověření}

\begin{itemize}
  \item \textbf{Fáze A} (čistě lineární algebra, $O(1)$ SVD + kombinatorika): $A\to$ null space
        $V\to$ Lemma 1/2 $\to$ všechny minimální cut-sety + $g_H$.
  \item \textbf{Fáze B} (čisté FEM, žádné RV): pro danou redukovanou konstrukci $b_{ij}$ (vliv
        zatížení), $a_{il}$ (vliv zbytkové síly) superpozicí.
  \item \textbf{Fáze C} (pravděpodobnost): (C.1)--(C.3) linearizovaná mezní funkce $\to$
        (C.4)--(C.6) Cornellův index sekvence (paralelní systém, \code{mvncdf}) $\to$ (C.3)
        nejhorší permutace cut-setu $\to$ (C.7)--(C.8) linearizace do ekvivalentní roviny $\to$
        (C.9) PNET kombinace nezávislých skupin $\to \beta_{\mathrm{sys}}$.
\end{itemize}

Celý pipeline byl ověřen na třech referenčních příkladech ze zdrojových článků (viz \code{metodika.md}):
3-prutová a 15-prutová příhrada (Wei \& Deng, shoda cut-setů), 3-podlažní příhrada
($\beta_{\mathrm{sys}}=3{,}0011$ implementace vs. $3{,}001162$ paper vs. $3{,}013374$ MCS $10^7$
vzorků -- rozdíl $-0{,}41\,\%$), a nezávisle přes brute-force progressive-collapse Monte Carlo
(\code{progressiveCollapseMCFn.m}, $\beta_{\mathrm{MC}}\approx3{,}02$--$3{,}04$).

\end{document}
